%! Rational Functions and End Behavior — Unit 3, Lesson 3.1 — Homework (Answer Key)
\documentclass[10pt]{article}
\usepackage{precalculus-article}
\usepackage{precalculus-key}   % pulls in -boxes; do NOT also load -boxes
\newcommand{\TallMath}[1]{$\displaystyle #1\rule[-1.4em]{0pt}{3.2em}$}

\begin{document}

\pageheader{Unit 3, Lesson 3.1}{Homework --- Answer Key}
\namedateperiod

\vspace{0.12in}

\begin{enumerate}[itemsep=14pt]

\item For each rational function below, identify the end behavior type (HA at $y = 0$,
HA at $y = k$, slant asymptote, or polynomial end behavior) and write both limit statements.

\medskip
\renewcommand{\arraystretch}{2.4}
\begin{tabular}{@{}p{3.8cm} p{3.2cm} p{3.0cm} p{3.0cm}@{}}
\textbf{Function} & \textbf{End Behavior Type} &
$\displaystyle\lim_{x\to\infty}$ & $\displaystyle\lim_{x\to -\infty}$ \\
\hline
$\dfrac{4x - 1}{2x + 3}$           & \ans{HA at $y=2$}     & \ans{$2$}       & \ans{$2$} \\
$\dfrac{x^2 + 1}{x^3 - 8}$         & \ans{HA at $y=0$}     & \ans{$0$}       & \ans{$0$} \\
$\dfrac{-6x^3}{2x^3 + x}$          & \ans{HA at $y=-3$}    & \ans{$-3$}      & \ans{$-3$} \\
$\dfrac{x^4 - x}{x^2 + 1}$         & \ans{Poly end behavior} & \ans{$+\infty$} & \ans{$+\infty$} \\
\end{tabular}

\item Explain why $\displaystyle\lim_{x \to \infty} \frac{7x^3 + 2}{x^3 - 1} = 7$.

\ansline{For large $x$, the term $7x^3$ dominates the numerator (the $+2$ is negligible),
and $x^3$ dominates the denominator (the $-1$ is negligible).
The quotient of leading terms is $7x^3/x^3 = 7$, so the outputs approach 7 as $x \to \infty$.}

\item A rational function $r$ satisfies $\displaystyle\lim_{x \to \infty} r(x) = 0$.
What can you conclude about the degrees of the numerator and denominator?

\ansline{The degree of the numerator must be strictly less than the degree of the denominator.
When the denominator dominates, the quotient of leading terms approaches 0, giving a HA at $y = 0$.}

\item A rational function has a slant asymptote.

\begin{enumerate}[label=(\alph*), itemsep=6pt]
  \item What must be true about the degrees of the numerator and denominator?

        \ansline{The degree of the numerator must equal the degree of the denominator plus 1
        (i.e., deg num $=$ deg denom $+ 1$).}

  \item Can you write the end behavior as a single limit value $\lim_{x\to\infty} r(x) = k$?

        \ansline{No. A slant asymptote means the quotient of leading terms is a non-constant
        linear expression, so the outputs grow without bound (not toward a constant $k$).
        The limit notation $\lim_{x\to\infty} r(x) = \pm\infty$ applies instead.}
\end{enumerate}

\item \textbf{AP-Style Multiple Choice.} Which function has a HA at $y = 3/2$?

\ans{(B)} $\dfrac{3x^2 - 1}{2x^2 + 7}$ --- the quotient of leading terms is $3x^2/2x^2 = 3/2$.

\textit{(A): deg num $<$ deg denom $\Rightarrow$ HA at $y = 0$.
(C): quotient is $2/3$.
(D): deg num $>$ deg denom $\Rightarrow$ polynomial end behavior.)}

\end{enumerate}

\vspace{0.1in}

\begin{reflectionbox}
\textbf{Lesson Preview.} In Lesson 1.8, we examine the \textit{zeros} of rational functions ---
the input values where the numerator equals zero.  Together with the asymptotes from today,
these zeros determine the overall shape of the graph of a rational function.
\end{reflectionbox}

\end{document}
